% ============================================================
%  sections/traction.tex
% ============================================================

\section{Traction \& Validation}

\subsection{Research foundation}

Runway is not a speculative idea. The core methodology is grounded in published
experimental research:

\begin{callout}
\textbf{Wahaj, S. (2026).} \textit{When Does Kubernetes Become Worth It? Measuring
Reliability, Stability, and Execution Overhead in Workflow Orchestration Systems
Moving from Shared VM to Isolated Kubernetes Execution.}
JENSEN Yrkeshögskola, DevOps24M programme.
\end{callout}

Key validated findings from 18 controlled experiments (108 individual pipeline runs):

\begin{itemize}
  \item VM reliability collapses non-linearly at 3 concurrent jobs (100\% → 66.7\% success rate)
  \item Kubernetes pod isolation maintains 100\% reliability across all tested concurrency levels
  \item Container startup overhead is predictable: 4.3s at L1 to 14.9s at L10, growing with concurrency
  \item The crossover point (VM → K8s migration threshold) is measurable and specific to hardware configuration
\end{itemize}

\textbf{This experimental dataset is Runway's training ground.} The recommendation rules in
v0.1 are calibrated against real failure data, not estimates.

\subsection{What validation is needed next}

\begin{enumerate}
  \item \textbf{5 pilot users} from the Dagster community Slack — confirm the latency trend
        signal is detectable and useful in real production pipelines.
  \item \textbf{3 documented predictions} — measure whether the breach date prediction is
        accurate within the stated confidence interval.
  \item \textbf{1 case study} — one team that acted on a Runway recommendation and avoided
        a production failure.
\end{enumerate}

These three items are achievable within 8 weeks of launch and constitute a fundable traction story.
